\subsubsection{ფრონტენდი}

ფრონტენდი React + TypeScript აპლიკაციაა (აწყობილი Vite-ით), რომლის
ცენტრშიც Monaco რედაქტორი დგას. მისი მთავარი ამოცანა ორმხრივი ნაკადის
გამართული მუშაობაა: ლოკალური keystrokes ბექენდისკენ, remote
ცვლილებები - რედაქტორისკენ.

\paragraph{ლოკალური ცვლილებების გაგზავნა.}
Monaco-ს ყოველი ცვლილება (\texttt{onDidChangeModelContent})
ითარგმნება IPC ბრძანებად: \texttt{insert}, \texttt{delete} ან
\texttt{replace}, პოზიციით და შიგთავსით. იმისთვის, რომ სწრაფი ბეჭდვისას
ბრძანებებმა ერთმანეთს არ გადაუსწრონ, ისინი ერთ ჯაჭვზე ეწყობა
(\texttt{opQueue}-ს მეშვეობით ყოველი შემდეგი ბრძანება წინას დასრულებას ელოდება).
თითო ბრძანებას თან მიჰყვება \texttt{baseSeq}, ბოლო remote ცვლილების
ნომერი, რომელიც რედაქტორს გაგზავნის მომენტში ჰქონდა ასახული, რომლითაც
ბექენდი პოზიციას გადათვლის როცა საჭიროა.

\paragraph{remote ცვლილებების მიღება.}
ბექენდის \texttt{crdt://remote-change} ივენთს ისმენს
\texttt{remoteChangeListener}, რომელიც ცვლილებას Monaco-ში პროგრამულად
ჩასვამს. აქ ორი ხაფანგი გვხვდება: ჯერ ერთი, პროგრამული ჩასმაც
\texttt{onDidChangeModelContent}-ს ააქტიურებს, რაც ცვლილებას უკან,
ბექენდისკენ გააბრუნებდა და უსასრულო ციკლს გააჩენდა, ამას ვკეტავთ
\texttt{isApplyingRemote} დროშით, რომელიც ჩასმის განმავლობაში ლოკალურ
დამმუშავებელს თიშავს. მეორე, წაკითხვის რეჟიმში მყოფ რედაქტორში Monaco
პროგრამულ ცვლილებებსაც უგულებელყოფს, ამიტომ remote ცვლილების ჩასმისას
\texttt{readOnly} დროშა წამიერად იხსნება. სნეფშოტის მიღებისას (მაგ.
სესიაში შესვლისას) ცალკე მსმენელი მთელ დოკუმენტს ანაცვლებს.

\paragraph{ინტერფეისის აგებულება.}
ვიზუალურად აპლიკაცია VS Code-ის სტილის გარსია: მარცხნივ ვიწრო 
side rail, რომელიც გვერდით პანელს ხსნის სამი სექციით -
ფაილები (გახსნა/შენახვა/ბოლო ფაილები), თანამშრომლობა (სესიის
წამოწყება/შეერთება/დატოვება) და პერსონალური პარამეტრები (სახელი, თემა,
ფონტის ზომა). ზედა ზოლში სესიის სტატუსი, მონაწილეთა ავატარები
(თითოეულს დეტერმინისტულად გამოთვლილი ფერი აქვს), მოსაწვევი ბმულების
popover და მენიუა. ინტერფეისის მდგომარეობა React ჰუკებშია
ორგანიზებული (\texttt{useRoomState}, \texttt{useSessionEvents},
\texttt{useWritePermission}, \texttt{useTheme} და ა.შ.), რომლებიც
ბექენდის ივენთებს უსმენენ და კომპონენტებს უბრალო props-ებით კვებავენ.

ფრონტენდის ვიზუალური დიზაინის იტერაციებში (ფერთა პალიტრა, კომპონენტების
განლაგება, თემები) დამხმარე ინსტრუმენტად ვიყენებდით Claude-ს
\cite{claude}; საბოლოო დიზაინის გადაწყვეტილებები და მათი იმპლემენტაცია
გუნდის წევრებისაა.

ფრონტენდს აქვს საკუთარი ტესტებიც (vitest): სინქრონიზაციის კრიტიკული
დამხმარე ლოგიკა, შებრუნებული რედაქტირებების გამოთვლა და
\texttt{baseSeq}-ის აღრიცხვა, ავტომატურად მოწმდება, ხოლო ტიპიზაცია და
ლინტერი CI-ში ეშვება.
